% id: 0304
% book: RPM 중학 3-1 · 03 근호를 포함한 식의 계산
% section: 유형 09 제곱근의 곱셈과 나눗셈의 도형에의 활용
% figure: crop:fig-0304.png
% answer: ④
% answer_source: 답지
% variant_level: 0 (원본 전사)
% 이 파일은 build.mjs 가 items.json 에서 만든다. 고칠 때는 items.json 을 고친다.
\smash{\raisebox{1.5mm}{\dmkichul{RPM 중학 3-1 0304번 · 44쪽 · 대표문제}}}
\noindent\begin{minipage}[t]{0.58\linewidth}\vspace{0pt}\raggedright 오른쪽 그림과 같이 직사각형~$\pt{ABCD}$에서 $\seg{AD}$, $\seg{CD}$를 각각 한 변으로 하는 정사각형을 그렸더니 그 넓이가 각각 $32$, $6$이 되었다. 이때 직사각형~$\pt{ABCD}$의 넓이는?\end{minipage}\hfill\begin{minipage}[t]{0.4\linewidth}\vspace{0pt}\centering \includegraphics[width=78.3pt]{figures/fig-0304.png}\end{minipage}\par
\begin{choices32}{\rule[-3ex]{0pt}{8ex}$4\sqrt{6}$}{\rule[-3ex]{0pt}{8ex}$6\sqrt{3}$}{\rule[-3ex]{0pt}{8ex}$7\sqrt{3}$}{\rule[-3ex]{0pt}{8ex}$8\sqrt{3}$}{\rule[-3ex]{0pt}{8ex}$6\sqrt{6}$}\end{choices32}
